\documentclass[11pt,a4paper]{article}
\usepackage[margin=25mm]{geometry}
\usepackage{amsmath,amssymb,amsthm,booktabs,hyperref}
\setlength{\emergencystretch}{2em}
\newtheorem{lemma}{Lemma}\newtheorem{theorem}[lemma]{Theorem}
\newcommand{\ON}{O'N}
\title{Excluding the type $(3,12)$ in the O'Nan group}
\author{Human Compression project}\date{}
\begin{document}\maketitle
\begin{abstract}
We prove that the O'Nan group is not generated by elements of orders
2 and 3 whose product has order 12. We count pairs according to the
subgroup they generate, rather than their possibly overlapping maximal
overgroups. Three proper subgroup families account for at least 954
of the 980 pairs with a fixed product. The remaining generating pairs
would occur in orbits of length 36, so there are none. The proof uses
the established O'Nan subgroup classifications, ordinary character data,
and the exhaustive catalogue of perfect groups of order 10752.
The finite calculations use only proper groups of order at most 3753792.
\end{abstract}

\section{Counting by the generated subgroup}
Throughout, $G=\ON$. We retain the established order
$|G|=460815505920$, simplicity of $G$, and its single conjugacy class
of elements of order 12. If $z$ is in that class, then $|C_G(z)|=36$.
The previously reviewed character calculation gives
\[
 |\{(x,z'):\ |x|=2,\ |z'|=12,\ |xz'|=3\}|
       =12544422105600.
\]
Dividing by $|G|/36$ therefore gives
\[
 T(z)=|\{x\in G:\ |x|=2,\ |xz|=3\}|=980.             \tag{1}
\]
The product-one character identity and its exact evaluation were proved
and reviewed in the common counting component. The named character
table, including the class and centralizer identifications, remains
an explicit external data premise.

For a finite subgroup $H\leq G$, put
\[
 e(H)=|\{(x,u)\in H^2:\ |x|=2,\ |u|=12,\ |xu|=3,
                         \ \langle x,u\rangle=H\}|.
\]
Pairs here are ordered. They determine the subgroup that they generate.
\begin{lemma}\label{partition}
The conjugacy class of $H$ contributes exactly
\[
                       \frac{36e(H)}{|N_G(H)|}           \tag{2}
\]
to $T(z)$, if we count only pairs generating a conjugate of $H$.
Contributions from distinct conjugacy classes of generated subgroups
are disjoint, even when those subgroups intersect nontrivially.
\end{lemma}
\begin{proof}
There are $|G|/|N_G(H)|$ conjugates of $H$. Each contains $e(H)$
pairs generating it; no such pair generates two different subgroups.
Thus the total is $|G|e(H)/|N_G(H)|$. This set of pairs is stable
under conjugation by $G$. Its second entries all have order 12, and
those elements form one $G$-class, so every possible second entry
has the same number of pairs. Divide by the class size $|G|/36$.
The same uniqueness of the generated subgroup proves disjointness.
\end{proof}
The formula can be used with a lower bound for $e(H)$ and an upper
bound for $|N_G(H)|$. It requires no multiplication of subgroup-cover
weights and no assumption about unique maximal overgroups.

\begin{lemma}\label{remainder}
The number of $x$ counted in $T(z)$ for which $\langle x,z\rangle=G$
is a multiple of 36.
\end{lemma}
\begin{proof}
The centralizer $C_G(z)$ acts on these $x$ by conjugation. A stabilizer
element centralizes both $x$ and $z$, hence centralizes the group they
generate. Simplicity makes the centre of $G$ trivial. Every stabilizer
is therefore trivial, and every orbit has size 36.
\end{proof}

\section{Two maximal classes and two small subgroups}
The complete maximal-subgroup classification gives two distinct
$G$-classes of $M\cong\operatorname{PSL}_3(7):2$, of order 3753792.
Both are self-normalizing: otherwise a maximal proper subgroup would
be normal in the simple group $G$.
The retained model is the action on the 57 points and 57 lines of the
projective plane over $\mathbb F_7$, extended by the graph involution.
The point action is projectivized matrix multiplication; the line
action uses the inverse transpose. These give the stated abstract
group. Its identification with the two maximal classes is borrowed
from the established classification, not inferred from its order.

The named table of $M$ has exactly one order-12 class, with centralizer
12. Thus the displayed element $u$ of order 12 satisfies $|C_M(u)|=12$;
no unverified class assignment for $u$ is needed. There are 132 retained, distinct involutions $x$ with
$|xu|=3$ and $\langle x,u\rangle=M$. For each of these pairs, an explicit word in $x,u$ has order 19.
The reader checks this from the literal cycles. The proper maximal
subgroups of $M$ have orders
\[
1876896,\ 8232,\ 1344,\ 672,\ 144,\ 144,\ 114.
\]
Only the first and last are divisible by 19. The first is the
index-two subgroup $L_3(7)$, preserving the two projective halves;
our displayed $u$ swaps those halves. The last has order 114,
which is not divisible by 12. Thus no proper maximal subgroup can
contain $x,u$, proving generation. This uses the established maximal
classification inside $M$, also recorded by rows 2 and 3 of the
subgroup roster.

Membership in $M$ is checked independently too. Given a retained
permutation preserving the halves, the images of the three coordinate
points and $[1,1,1]$ determine a candidate projective matrix up to scalar.
The reader solves this linear system over $\mathbb F_7$, checks all
57 point images and all 57 inverse-transpose line images, and checks
that the determinant is a cube. This last condition is exactly
membership in $\operatorname{PSL}_3(7)$ inside
$\operatorname{PGL}_3(7)$. For a permutation swapping the halves,
the reader first composes with the fixed point--line interchange.
The retained coordinates are checked against the complete projective
point set. No group-order computation is used to certify the 132 pairs.

Conjugating to the $|M|/12$ possible second entries already gives
\[
                        e(M)\geq 132|M|/12=11|M|.       \tag{3}
\]
We need only this lower bound. The retained computation also found ten
additional pairs generating smaller subgroups, but none is included
in (3). By (2), each maximal class contributes at least 396, and the
two classes contribute at least 792.

We next use the maximal subgroup $B$ described as $4^3.L_3(2)$.
Here the classification asserts a normal subgroup $C_4^3$ with
quotient $L_3(2)$; in particular $|B|=10752$ and $O_2(B)=C_4^3$.
The exact table \texttt{4\^{}3.L3(2)} has one linear character.
Since the number of linear characters is $|B/B'|$, it follows that
$B$ is perfect. This table identification is another explicit named
data premise.

We identify the finite model of $B$ as follows. The perfect-group
catalogue provides all isomorphism types of order 10752. It contains
nine types at this order. The retained calculation of their largest
normal 2-subgroups leaves precisely catalogue entries 3 and 4 with
$O_2=C_4^3$; in all seven other entries that subgroup has exponent 2.
The catalogue's exhaustiveness is a classification premise, stated in
the GAP Reference Manual, Section 50.6, and derived from the work of
Holt and Plesken. The historical omitted orders listed there do not
include 10752.

Both remaining models have two classes of elements of order 12,
each with centralizer 12. In entry 3 there are 12 involutions $x$
with $|xu|=3$ at either such class; in entry 4 there are 36.
For the named table of $B$, the class multiplication coefficients
give 36 at both classes. Thus $B$ cannot be entry 3 and is isomorphic
to entry 4. This is exclusion within an exhaustive list of
isomorphism types; matching a numerical fingerprint without that
classification would not identify the group.

The retained model of entry 4 contains explicitly generated subgroups
$H_{48}$ and $H_{768}$ of orders 48 and 768. The independent literal
reader obtains the following data.
\begin{center}\begin{tabular}{rrrrr}\toprule
$|H|$&involutions&12-classes&size of each class&generating $x$ per $u$\\\midrule
48&7&4&4&3\\
768&55&4&64&24\\\bottomrule
\end{tabular}\end{center}
Consequently
\[
                 e(H_{48})=4\cdot4\cdot3=48,
 \qquad e(H_{768})=4\cdot64\cdot24=6144.               \tag{4}
\]
These are counts of pairs generating the small group itself, not all
pairs lying in it. The reader enumerates the groups from their literal
permutations, checks closure, forms the conjugacy orbits, and tests the
subgroup generated by each eligible pair. It also checks that their
generators belong to entry 4. Transport by the proved isomorphism
$B\cong\text{entry 4}$ gives subgroups of $G$ with exactly these
intrinsic counts. Their particular $G$-class labels are unnecessary.

\section{Normalizer bounds from the subgroup roster}
We use the previously acquired complete 581-class subgroup roster
of $G$ as an attributed classification premise. Its class lengths
are numbers of conjugate subgroups, so a row of length $c$ has
normalizer order $|G|/c$. The four rows of order 768 are 97--100.
Each has class length 300010095, and hence normalizer order 1536.
It follows that
\[
                         |N_G(H_{768})|=1536.            \tag{5}
\]
The rows of order 48 are 382--395. Apart from row 395 their
normalizers have order 48 or 96. Row 395 has normalizer order 384,
but cannot be the class of $H_{48}$: its maximal subgroups are
row 528, of order 16, and row 545, isomorphic to $A_4$.
Neither can contain a cyclic group of order 12. Any such cyclic
subgroup of row 395 would lie in a maximal proper subgroup, a
contradiction. Since $H_{48}$ contains its displayed element $u$
of order 12, we have
\[
                         |N_G(H_{48})|\leq96.            \tag{6}
\]
Only the two indicated maximal-subgroup entries are used in this
exclusion. Cover multiplicities are not multiplied or treated as a
full incidence matrix. The certificate joins the numerical rows and
the two subgroup labels directly to native roster pages 5, 14, 15,
20 and 21.

\section{The zero and its source scope}
\begin{theorem}
There are no generators $r,s$ of $G$ with $|r|=3$, $|s|=12$ and
$|rs|=2$.
\end{theorem}
\begin{proof}
The generated subgroup orders in the three families are respectively
3753792, 768 and 48, so the families are disjoint. Equations (2)--(6)
give at least
\[
  792+\frac{36\cdot6144}{1536}+\frac{36\cdot48}{96}
                       =792+144+18=954                 \tag{7}
\]
proper pairs at each fixed element $z$ of order 12. By (1), at most
$980-954=26$ pairs remain available to generate $G$. Lemma~\ref{remainder}
makes this a nonnegative multiple of 36; it must be zero.

If $r,s$ were the claimed generators, put $x=rs$ and $z=s^{-1}$.
Then $|x|=2$, $|z|=12$, $xz=r$ has order 3, and
$\langle x,z\rangle=\langle r,s\rangle=G$, contrary to the zero
just proved. Conversely these substitutions recover $r=xz$ and
$s=z^{-1}$, so the conventions agree exactly.
\end{proof}

This supplies the remaining zero of type $(3,12)$ in the fixed L046
source. It proves the conclusion attributed there to Al-Kadhi, without
using the incomplete inference from a single proper pair identified in
the earlier source audit. We do not attribute that inference to the
present proof. Subject to distinct review, exact table coverage advances
from 55 to 56 of 136 cells; the other 80 cells remain separate.
The positive chirality witnesses and the exceptional type $(3,15)$ are
unaffected. This is a proof of an existing asserted zero, not a new
conjecture settlement or a complete treatment of the original paper.

The maximal and subgroup classifications, named character tables,
and perfect-group catalogue are borrowed. The contribution is the
explicit proper-subgroup counts and their combination with the free
centralizer action. The working record discloses the collaborative
derivation of (7); the final mathematical review is by a different
agent from its two authors. The computational checks do not constitute
a Lean formalization.
\begin{thebibliography}{9}
\bibitem{CL} T. Connor and D. Leemans, \emph{Algorithmic enumeration of
regular maps}, Ars Math. Contemp. 10 (2016), 211--222.
\bibitem{Wilson} R. A. Wilson, \emph{The maximal subgroups of the O'Nan
group}, Journal of Algebra 97 (1985), 467--473.
\bibitem{Atlas} Atlas of Group Representations, retained O'Nan page;
the acquired O'Nan subgroup roster of Connor and Leemans.
\bibitem{HP} D. F. Holt and W. Plesken, \emph{Perfect Groups},
Clarendon Press, Oxford, 1989; GAP Reference Manual, Section 50.6.
\end{thebibliography}
\end{document}
